% Related Work
% Organize by themes/categories, not paper-by-paper
% Show how your work builds on, differs from, or fills gaps in existing work

\label{sec:related}

% Introduction to related work
Wikipedia’s unique open editing model offers an opportunity to research its governance in relation to politics and free speech laws.

% Subsection 1: First category of related work
\subsection{Wikipedia Governance and Community Culture}

Reagle's seminal work on Wikipedia culture~\cite{reagle2010good} examines the collaborative practices and governance mechanisms that enable Wikipedia's success. 

% Subsection 2: Second category of related work
\subsection{Wikipedia Trust and Reputation}
de Laat studied how Wikipedia monitors edits in order to avoid vandalism, with both admins and autonomous bots~\cite{de2014open}. Their paper states that Wikipedia experiences roughly 9000 vandalism edits daily, which is 7\% of all edits. It elaborates on institutional trust, which in this context is how Wikipedia approaches trusting contributors. Wikipedia uses backgrounding trust, which requires monitoring reputation and uses corrective mechanisms in the background to enforce norms, with sanctions placed on users who deviate. Bots are used to highlight edits that are potentially vandalism, which raises the question of bias within the automated bots. 

A paper on chilling effects showed that awareness of surveillance of IP address impacted Wikipedia usage, which is relevant because it means that potential editors/admins in more repressive countries will not participate to the full extent due to worries of surveillance~\cite{penney2016chilling}. The paper uses empirical backing to show the cost of spam control, with limits on edits, potentially especially for politically controversial pages.

\subsection{Wikipedia and Political Context}
The study done by Parra~\cite{parra2024righting} uses a casual inference framework and showed that politicians who switched from left to right wing parties had a drop in sentiment on their Wikipedia articles, which could show systematic political bias. This means that even with Wikipedia's monitoring of vandalism there is bias, making it important for us to study the role of free speech and politics with Wikipedia users. An Oxford paper found that the majority of articles about places were edited by people living in the UK, USA, France, Germany, and Italy~\cite{Oxford2015WikipediaWorldView}. Digital connectivity is a contributing factor, along with high income levels.

% Positioning your work
\subsection{Our Work in Context}

The majority of prior existing similar research examines survey-based measures; we are using quantitative empirical behavioral data to examine user participation on Wikipedia. Our research is primarily focused on observable edits and behavior. Additionally, we focus on several different countries and create a bridge between quantitative edit evidence and free speech laws related to those countries. This contributes to further understanding how free speech laws can impact user behavior in digital public squares like Wikipedia.
